\chapterimage{apps.jpg}

\chapter{Capa 7: Aplicació}\label{sec:layer7}

\begin{minipage}{0.4\linewidth}
\begin{center}
\begin{bytefield}{16}
\bitbox{16}{\color{color1} Layer 7: Application} \\
\bitbox{16}{Layer 4: Transport} \\
\bitbox{16}{Layer 3: Internet} \\
\bitbox{16}{Layer 2: Network (LAN)} \\
\bitbox{16}{Layer 1: Physical} \\
\end{bytefield}
\end{center}
\end{minipage}
\begin{minipage}{0.6\linewidth}
\begin{center}
\includegraphics[width=\linewidth]{network_layer7.pdf}
\end{center}
\end{minipage}

La capa superior de la pila de xarxa TCP/IP inclou totes les aplicacions
que utilitzen TCP o UDP per al transport. És, amb diferència, la que conté
el nombre més gran de protocols, i on aquests protocols es creen i se substitueixen
més ràpidament.

A més de les aplicacions d'usuari, hi ha diversos sistemes i protocols auxiliars que
ajuden en la configuració i el funcionament de tots els hosts (equips connectats a la xarxa) d'Internet.
En aquest capítol s'aborden alguns dels més crítics.

\section{DNS -- Sistema de noms de domini}\label{sec:layer7:dns}

El protocol \concept{DNS} és responsable de traduir \conceptRef{domain name}{noms de domini}
com ara \otherBase{ietf.org} o\\\otherBase{cv.uab.cat} a adreces IP
que es poden posar dins dels datagrames.

Els serveis DNS funcionen tant sobre TCP com sobre UDP, \ie, es poden habilitar
tots dos protocols de transport. En tots dos casos, s'utilitza el port \otherBase{53}.

\subsection{Jerarquia}\label{sec:layer6:dns:hierarchy}

\begin{center}
\includegraphics[width=0.65\linewidth]{dns_hierarchy.pdf}
\end{center}

Els noms de domini d'Internet són seqüències de noms separats per punts.
Aquests noms es presenten del més específic al més general de manera jeràrquica.
Per exemple, a \otherBase{cv.uab.cat}, \otherBase{cv} pertany a \otherBase{uab.cat},
i \otherBase{uab} a \otherBase{cat}.
% 
El nom de més a la dreta, \eg, \otherBase{cat} o \otherBase{com},
ha de ser un dels \conceptRef{top-level domain}{dominis de primer nivell} de la IANA.

Els servidors DNS s'organitzen de la mateixa manera jeràrquica que aquests noms de domini.
Hi ha $13$ servidors \concept{root DNS} (de \inlineCode{Root-A} a \inlineCode{Root-M}),
que són responsables de tots els dominis de primer nivell.
Pots preguntar a qualsevol d'ells per \otherBase{com},
però també per \otherBase{su}, i en sabran la resposta.

Això no vol dir que cada servidor arrel conegui tots els dominis \textit{dins}
del seu domini de primer nivell (\eg, \otherBase{uab.cat}): només coneixen aquests
\href{https://en.wikipedia.org/wiki/List_of_Internet_top-level_domains}{\underline{dominis de primer nivell}}.
(\eg, \otherBase{cat}).
% 
La feina dels servidors DNS arrel és redirigir-te a un altre servidor DNS encarregat
del \conceptRef{top-level domain}{domini de primer nivell} que has demanat.
 
De manera semblant, els servidors DNS responsables de \otherBase{cat} coneixeran
qualsevol domini \otherBase{*.cat} existent (\eg, \otherBase{uab.cat}),
però no els subdominis (\eg, \otherBase{cv.uab.cat}).
% 
La feina d'aquests servidors DNS de domini de primer nivell és redirigir-te a
altres servidors DNS \conceptRef{authoritative}{autoritatius} per al domini de segon nivell
sol·licitat (\eg, \otherBase{uab.cat}).
% 
Els dominis de segon nivell poden llavors estructurar la seva \conceptRef{zone}{zona} de noms
com vulguin, incloent-hi qualsevol nombre i profunditat de subdominis, així com
la traducció de nom a IP \textit{dins} de la zona d'autoritat.

\begin{exercise}
Tothom fa servir el DNS constantment. Quins serien els riscos de tenir
un sistema de traducció de noms centralitzat en lloc de la naturalesa distribuïda del DNS?
\end{exercise}

\subsection{Registres}\label{sec:layer7:dns:records}

El DNS es pot veure com una base de dades distribuïda amb diferents tipus d'entrades
anomenades \conceptRef{DNS record}{registres DNS}. Els principals són:

\begin{itemize}
\item[\textbf{A}:] Un \conceptRef{A}{registre A} associa un nom de domini a una adreça IPv4.
\item[\textbf{AAAA}:] Un \conceptRef{AAAA}{registre AAAA} associa un nom de domini a una adreça IPv6.
\item[\textbf{NS}:] Un \conceptRef{NS}{registre NS} apunta a un servidor DNS que pot ajudar.
\item[\textbf{MX}:] Un \conceptRef{MX}{registre MX} proporciona el nom de domini del servidor de correu electrònic
  que gestiona els missatges d'un domini.
\item[\textbf{CNAME}:] Un \conceptRef{CNAME}{registre CNAME} associa un nom de domini a un altre nom de domini.
\end{itemize}

Les aplicacions fan \conceptRef{request}{peticions} als servidors DNS indicant el tipus
de registres que els interessen. Les \conceptRef{reply}{respostes} del servidor poden ser
els registres A o MX demanats,
però el client també pot rebre una redirecció a un altre DNS (NS)
o un àlies a un altre nom de domini (CNAME).

\subsection{Resolució}\label{sec:layer7:dns:resolution}

La naturalesa distribuïda i jeràrquica de la base de dades del DNS fa que
de vegades calgui intercanviar múltiples missatges per satisfer una \conceptRef{query}{consulta} DNS.
 
Per exemple, l'Alice pot fer una \conceptRef{iterative}{resolució iterativa} completament tota sola.
Les dues primeres respostes contenen un registre NS que apunta al següent
servidor DNS i l'adreça IP d'aquest servidor DNS, però no l'adreça IP del
\otherBase{cv.uab.cat} demanat.

\begin{center}
\includegraphics[width=\linewidth]{dns_iterative.eps}
\end{center}

El proveïdor d'Internet (\concept{ISP}) o l'empresa de l'Alice probablement oferiran servidors DNS
la feina dels quals és rebre les peticions de l'Alice i resoldre-les per ella de manera
recursiva.

\begin{center}
\includegraphics[width=1.05\linewidth]{dns_recursive.eps}
\end{center}

La combinació de la \conceptRef{recursive}{resolució recursiva} i les taules de \conceptRef{cache}{memòria cau} fa
del DNS un sistema força eficient, amb una disponibilitat altíssima i una latència baixa.

\begin{exercise}
Pots utilitzar les eines \inlineCode{host} i \inlineCode{dig} per fer consultes DNS
i veure'n els resultats. Per exemple, la meva sortida de \inlineCode{dig uab.cat -t ns}
conté les línies següents:

\vspace{0.25cm}
\begin{verbatim}
;; QUESTION SECTION:
;uab.cat.                       IN      NS
;; ANSWER SECTION:
uab.cat.                172800  IN      NS      dns.uab.es.
;; ADDITIONAL SECTION:
dns.uab.es.             172112  IN      A       158.109.0.1
\end{verbatim}
\vspace{0.25cm}

\begin{itemize}
\item Quin tipus de registre he demanat?
\item Quins dos tipus de registre he rebut?
\item Per què he rebut dos tipus de registre?
\item Fes una resolució iterativa manual de \otherBase{cv.uab.cat}.
  Pots començar amb alguna cosa com \inlineCode{dig @202.12.27.33 cv.uab.cat}
  per consultar el servidor arrel M.
\end{itemize}
\end{exercise}


\section{DHCP -- Protocol de configuració dinàmica de hosts}\label{sec:layer7:dhcp}

Per poder participar plenament a Internet, un host ha de tenir
correctament configurats els aspectes següents:
\begin{itemize}
\item \textbf{Adreça IP}:
L'adreça IP és necessària per identificar un node a la xarxa, i perquè aquest node sàpiga quins missatges van adreçats a ell.
\item \textbf{Màscara de subxarxa}:
Els nodes l'han de conèixer per determinar quan han de lliurar els datagrames directament, i quan indirectament.
\item \textbf{IP del gateway (passarel·la)}:
Cal l'adreça IP d'almenys un router (encaminador) per lliurar indirectament missatges amb destinació fora de la LAN (xarxa d'àrea local) actual.
\item \textbf{Adreça IP del servidor DNS}:
Tot i que no és estrictament necessària per utilitzar TCP/IP, l'adreça d'un servidor DNS és extremadament útil, i necessària per resoldre
noms de domini a adreces IP.
\end{itemize}

Tradicionalment, aquests paràmetres s'havien de configurar manualment a cada ordinador.
Avui dia, la majoria dels ordinadors d'usuari es configuren automàticament amb \concept{DHCP}.
Els sistemes operatius moderns contenen un client DHCP que consulta periòdicament la LAN fins que queda configurat.

\begin{remark}
El maquinari de xarxa ve amb una adreça MAC per defecte preconfigurada.
A més, un node necessita conèixer la seva adreça MAC per fer funcionar el protocol DHCP.
Per tant, les adreces MAC no es configuren mitjançant DHCP.
\end{remark}

\newpage
DHCP funciona sobre UDP. Els servidors DHCP escolten al port \otherBase{67}, i els clients
escolten al port \otherBase{68}. L'intercanvi DHCP ideal consta de $4$ missatges:
Discovery, Offer, Request, Acknowledgement.

\begin{center}
\includegraphics[width=\linewidth]{dhcp_dora.eps}
\end{center}

\begin{itemize}

\item \textbf{Discovery}\\
L'objectiu d'aquest missatge és trobar servidors DHCP disposats a proporcionar
una configuració.\\[-0.3cm]

En la primera comunicació, el client
no coneix la IP del servidor DHCP. En lloc d'això, utilitza:
\begin{itemize}
\item La seva adreça MAC com a adreça MAC d'origen.
\item L'adreça de broadcast (difusió) de la LAN (\otherBase{ff:ff:ff:ff:ff:ff}) com a adreça MAC de destinació.
\item \otherBase{0.0.0.0} com a IP d'origen, que vol dir ``desconeguda''.
\item L'adreça de \concept{local broadcast}, \otherBase{255.255.255.255}, com a IP de destinació.
\end{itemize}

\vspace{0.25cm}
\item \textbf{Offer}\\
Zero o més servidors respondran al client amb ofertes de configuració.
Aquestes configuracions (\eg, adreces IP) no són propietat del client
fins que el servidor envia un DHCP Acknowledgement.\\[-0.3cm]

El servidor porta el control d'un \conceptRef{pool}{pool} (conjunt) d'adreces que
assigna dinàmicament als clients.
% 
Les adreces IP s'ofereixen durant un període de temps fix, després del qual caduquen
tret que es \conceptRef{renew}{renovin} (vegeu més avall). Les adreces que no es renoven tornen
al conjunt i s'ofereixen de nou als clients a mesura que fan peticions.
\\[-0.3cm]

\begin{itemize}
\item El servidor pot utilitzar la seva adreça IP real i la que ofereix al client.
\item El client sap que n'és el destinatari perquè el servidor té en compte
l'adreça MAC d'origen del Discovery del client.
\end{itemize}

\vspace{0.25cm}
\item \textbf{Request}\\
El client tria una de les ofertes rebudes i demana que
es completi l'assignació.\\[-0.3cm]

\begin{itemize}
\item El client utilitza l'adreça de broadcast local, \otherBase{255.255.255.255}, com a \mbox{destinació}.
\item El client manté \otherBase{0.0.0.0} com a origen perquè l'adreça oferta encara no és seva.
\end{itemize}

\vspace{0.25cm}
\item \textbf{Acknowledgement (ACK)}\\
El servidor confirma la petició. Un cop el client rep aquest ACK,
comença a utilitzar la configuració oferta com a pròpia.
Aquest ACK no s'ha de confondre amb el de TCP.
\\[-0.3cm]

\vspace{0.25cm}
\item \textbf{Negative Acknowledgement (NACK)}\\
També és possible que el servidor denegui la petició, \eg, si
ha passat massa temps des de l'oferta. En aquest cas,
s'envia un Acknowledgement negatiu (NACK). Si es rep un NACK,
el client ha de tornar a la fase de Discover.
\end{itemize}

Quan un servidor DHCP \conceptRef{lease}{cedeix} una adreça a un client, ho fa
durant un període de temps limitat $T$. Aquest període el fixa el servidor DHCP, i es pot configurar.
Un cop transcorregut $T/2$, els clients poden intentar renovar la configuració cedida. Per fer-ho,
es repeteix el cicle Request-Acknowledgement al llarg del temps fins que un dels dos extrems no el completa a temps,
o bé el servidor DHCP envia un NACK.

\begin{center}
\includegraphics[width=\linewidth]{dhcp_renewal.eps}
\end{center}

\begin{exercise}
El router domèstic típic implementa alhora un client DHCP i un servidor DHCP.
Explica com és possible i per què és útil.
\end{exercise}


\section{Sistemes autònoms}\label{sec:layer7:as}

L'encaminament es fa \conceptRef{hop}{salt} a salt, utilitzant \conceptRef{routing table}{taules d'encaminament}
tal com s'ha explicat a la Secció~\ref{sec:layer3:routing}. Aquestes taules s'han de coordinar a gran
escala per garantir la connectivitat del graf amb tots els hosts d'Internet.

\begin{center}
  \vspace{-0.2cm}
 \includegraphics[width=\linewidth]{as.pdf}
 \vspace{-0.75cm}
\end{center}


\subsection{Organització}\label{sec:layer7:as:organization}

Per facilitar aquesta coordinació, les LAN s'agrupen en \conceptRef{AS}{sistemes autònoms} (\concept{AS}).
Aquests gestionen blocs d'adreces IP que no se solapen (\otherBase{w.x.y.z/N}), amb propietat i gestió separades.
Moltes empreses, universitats i ISP tenen els seus propis AS, que estan
\href{https://2ip.io/analytics/asn-list/}{\underline{llistats}} públicament. Per exemple, la xarxa de la UAB pertany a
l'AS13041.

Els \conceptRef{IXP}{punts d'intercanvi d'Internet} (\concept{IXP}) són un tipus especial d'AS el propòsit dels quals és interconnectar altres AS,
en lloc d'allotjar serveis i/o clients. Per descomptat, fer funcionar aquests IXP no és gratuït, així que s'estableixen relacions
comercials amb els AS interconnectats.
%
Alguns IXP es comprometen a donar servei a qualsevol que
compleixi els seus requisits bàsics, no estan controlats per cap operador únic, i ofereixen preus transparents i sovint equitatius
als seus clients. Aquest tipus d'IXP s'anomena \conceptRef{NP}{punt neutre} (\concept{NP}), en contraposició
als punts no neutres utilitzats, \textit{p.\,ex.}, per un sol operador de xarxa.

\newpage
Tots els AS fan funcionar \conceptRef{BR}{routers frontera} (\concept{BR}) per interconnectar AS. Quan un router frontera rep un datagrama
amb destinació dins de l'AS, el reenvia a través dels \conceptRef{internal router}{routers interns} (IR) de l'AS.
De manera semblant, quan un host dins d'un AS necessita contactar amb un host d'un altre AS, els datagrames es reenvien a través de routers frontera
fins que arriben a l'AS de destinació i, finalment, a la LAN de destinació.

\subsection{Configuració de la taula d'encaminament}\label{sec:layer7:as:routing_table_configurati}

Configurar manualment les \conceptRef{routing table}{taules d'encaminament} de tots els routers frontera i routers interns
seria massa costós i propens a errors. En lloc d'això, s'utilitzen diferents algorismes per configurar-les automàticament.
\begin{itemize}
\item Dins de l'AS, els routers interns i els routers frontera utilitzen un dels
\conceptRef{interior gateway protocol}{protocols de gateway interior} disponibles.
Sovint s'utilitza \concept{RIP} (Routing Information Protocol) per la seva configuració relativament senzilla.
Dit això, els AS poden triar lliurement l'algorisme que utilitzen internament.\\[-0.2cm]

RIP implementa una variant de l'algorisme de camí més curt de Bellman-Ford per minimitzar el nombre de \conceptRef{hop}{salts}
que els datagrames han de fer dins de l'AS. Periòdicament (\eg, cada $30$~s), cada IR i BR:\\[-0.4cm]
% 
  \begin{enumerate}
  \item Anuncia la seva distància \textit{percebuda} als altres routers de l'AS.
  \item Actualitza les seves distàncies percebudes a partir dels anuncis dels altres routers.
  \item Actualitza en conseqüència la seva \conceptRef{routing table}{taula d'encaminament} cap a totes les LAN de l'AS.
  \end{enumerate}

\begin{exercise}
RIP utilitza el nombre de salts com a mètrica de cost.
% 
Per què podria ser problemàtic?
% 
Com es pot resoldre aquest problema?
\end{exercise}

\vspace{0.25cm}
\item Entre AS, només cal configurar els routers frontera; els routers interns simplement han de dirigir
el trànsit exterior a través dels routers frontera. \\[-0.3cm]

Els routers frontera executen un \conceptRef{exterior gateway protocol}{protocol de gateway exterior} a més del
protocol de gateway interior utilitzat dins de l'AS. Els diferents AS s'han de posar d'acord en
l'algorisme que utilitzen, de manera que normalment s'utilitza \concept{BGP} (border gateway protocol).\\[-0.3cm]

\concept{BGP} és un protocol més complex que \concept{RIP}. Aquesta complexitat
és necessària perquè BGP coordina routers frontera que pertanyen a AS diferents, \ie,
amb propietaris diferents:\\[-0.4cm]
\begin{itemize}
\item RIP utilitza un únic nombre (la distància, \ie, el nombre de \conceptRef{hop}{salts})
per guiar un algorisme de camí més curt sobre un graf. BGP utilitza \conceptRef{path}{camins} en lloc de distàncies.
Això permet a BGP triar rutes amb més \conceptRef{hop}{salts}, però més convenients per a l'AS,
\eg, en termes de temps, cost econòmic, càrrega interna, etc.\\[-0.4cm]

\item BGP es va dissenyar tenint en compte la seguretat i la confiança, i inclou mecanismes
per verificar l'autenticitat dels camins anunciats per un router frontera.
\end{itemize}
\end{itemize}

\begin{remark}
BGP també es pot utilitzar com a protocol de gateway interior de l'AS.
Normalment això només ho fan AS molt grans.
\end{remark}

\begin{exercise}
Com pot BGP traduir interessos econòmics en configuracions d'encaminament?
\end{exercise}

